\documentclass[AutoFakeBold,AutoFakeSlant]{beamer}
\usepackage{ctex}
\usepackage{hyperref}
\usepackage[defaultsans]{droidsans}
\usepackage{fontspec}
\usepackage{newtxmath}
\usepackage{latexsym,amsmath,xcolor,multicol,booktabs,calligra,tcolorbox}
\usepackage{graphicx,tabularx,multirow,makecell}
\usepackage{listings}
\usepackage{tikz}
\usetikzlibrary{arrows.meta,positioning}
\lstset{basicstyle=\ttfamily\tiny,breaklines=true,breakatwhitespace=true,columns=fullflexible,frame=single,backgroundcolor=\color{gray!6}}
\usetheme{Madrid}
\usecolortheme{default}
\setbeamertemplate{navigation symbols}{}
\setbeamertemplate{footline}[frame number]
\graphicspath{{./figures_part1/}{../src/analyzer/output/run_20260427_001723/}}

\author{李禹江，罗嘉恒}
\title{礼物的流动：基于 Multiagent 的礼物流动现象社会学模拟仿真实验}
\subtitle{受阎云翔《礼物的流动》启发的 100 人村庄 365 天长周期仿真}
\institute{选题四：Multi-agent}
\date{2026-07-26}

\begin{document}
\kaishu

\begin{frame}
    \titlepage
\end{frame}

\begin{frame}
    \frametitle{目录}
    \tableofcontents
\end{frame}

\section{研究背景与问题}

\begin{frame}
    \frametitle{理论锚点：阎云翔《礼物的流动》}
    \small
    \begin{itemize}
        \item \textbf{核心命题}：礼物交换并非纯经济行为，而是嵌入道德、权力与时间性的社会过程。
        \item \textbf{四条可检验民族志模式}：
        \begin{enumerate}
            \item 延迟互惠——欠情与还情之间存在显著时间差；
            \item 纵向不对称——弱关系向高地位者的"进贡式"礼物不进入一对一回礼窗口；
            \item 互惠中态——多数偿还、少量违约、非即时对等；
            \item 声望吸引——礼物流向高声望节点集中，而非纯财富驱动。
        \end{enumerate}
        \item \textbf{方法主张}：不是"让 LLM 聊天看起来像人"，而是\textbf{显式社会规则 + 个体认知状态 + 主观记忆 + 可选话语}的混合多智能体，可复现民族志式宏观中态，并保留可拆解的微观机制。
    \end{itemize}
\end{frame}

\begin{frame}
    \frametitle{研究问题与可检验假说}
    \small
    \begin{columns}
        \begin{column}{0.48\textwidth}
            \textbf{研究问题}
            \begin{enumerate}
                \item \textbf{RQ1} 礼物流是否向高声望节点集中？互惠/违约是否中态？
                \item \textbf{RQ2} 回礼窗口、关系轴、经济约束、认知驱动是否各自解释得通？
                \item \textbf{RQ3} 回礼是否延迟互惠？决策/话语是否与压力、欠情、讨厌一致？
                \item \textbf{RQ4} 关系/性格/社会属性/经济状况在什么程度上影响集体事务解决？
            \end{enumerate}
        \end{column}
        \begin{column}{0.48\textwidth}
            \textbf{可检验假说}
            \begin{itemize}
                \item \textbf{H1} 声望吸引：$r_{\text{prestige}\times\text{gift}}>0$
                \item \textbf{H2} 互惠中态：reciprocity 0.7–0.95，default 0.02–0.20
                \item \textbf{H3} 纵向不对称：进贡不进债窗
                \item \textbf{H4} 延迟互惠：delay $>0$
                \item \textbf{H5} 状态—行为一致：话语与情绪/礼债一致
            \end{itemize}
        \end{column}
    \end{columns}
\end{frame}

\section{实验框架：9 个核心组件}

\begin{frame}
    \frametitle{核心链路与 9 组件总览}
    \centering
    \begin{tikzpicture}[
        node distance=4mm and 5mm,
        box/.style={draw, rounded corners, align=center, minimum width=2.0cm, minimum height=6mm, font=\scriptsize},
        arr/.style={-{Latex[length=1.6mm]}, thick},
    ]
        \node[box, fill=blue!10] (event) {8 事件系统\\公共/社交/回礼};
        \node[box, fill=green!14, right=of event] (dlg) {6 LLM 对话\\受状态叙事约束};
        \node[box, fill=cyan!12, right=of dlg] (pet) {5 PET\\个人事件表};
        \node[box, fill=orange!16, below=of pet] (ct) {4 认知树\\主观关系记忆};
        \node[box, fill=violet!10, left=of ct] (bdie) {3 BDIE\\信念/欲望/意图/情绪};
        \node[box, fill=yellow!18, left=of bdie] (pers) {9 性格派生\\BigFive};
        \node[box, fill=red!12, below=of bdie] (rule) {1 RuleEngine+礼单\\唯一写世界};
        \node[box, fill=gray!14, left=of rule] (rel) {2 关系图\\轴/亲缘/讨厌};
        \node[box, fill=blue!8, right=of rule] (eco) {7 经济/日活力\\生存约束};
        \draw[arr] (event) -- (dlg);
        \draw[arr] (dlg) -- (pet);
        \draw[arr] (pet) -- (ct);
        \draw[arr] (ct) -- (bdie);
        \draw[arr] (bdie) -- (rule);
        \draw[arr] (pers) -- (bdie);
        \draw[arr] (rel) -- (rule);
        \draw[arr] (eco) -- (rule);
    \end{tikzpicture}
    \vspace{0.3em}
    \scriptsize
    核心链路：事务触发 $\to$ 对话叙事 $\to$ RuleEngine 改关系/经济/礼物 $\to$ JSONL/快照 $\to$ 像素步进回放
\end{frame}

\begin{frame}
    \frametitle{9 组件分工总表}
    \tiny
    \begin{tabularx}{\textwidth}{|c|l|X|X|}
        \toprule
        \textbf{\#} & \textbf{组件} & \textbf{模拟人类的哪一块} & \textbf{关掉会怎样} \\
        \midrule
        1 & RuleEngine + 礼单账本 & 人情规范与结算（送/还/违约/场合） & 互惠崩溃或账目不可信 \\
        \midrule
        2 & 关系图（轴/亲缘/讨厌） & 社会结构与情感取向 & H3 失效；闲话无靶 \\
        \midrule
        3 & BDIE（信念/欲望/意图/情绪） & 当下身心状态与行动优先级 & 社交与回礼变"平均人" \\
        \midrule
        4 & 认知树 & 主观记忆与关系叙事 & 对话失语境；单向流动增强 \\
        \midrule
        5 & PET 个人事件表 & 近期显著事件记忆 & 事件提及与针对性回应变弱 \\
        \midrule
        6 & LLM 对话系统 & 自然语言表达与局部态度 & 只剩指标，缺可演示微观互动 \\
        \midrule
        7 & 经济系统 / 日活力 & 生存约束与压力底噪 & 礼期望膨胀、分层扭曲 \\
        \midrule
        8 & 事件系统 & 事务触发（公共/社交/回礼） & 时序失真，burstiness 消失 \\
        \midrule
        9 & 性格 / 人格派生 & BigFive 派生（内向/挑剔/记仇） & 决策同质化，微观异质性消失 \\
        \bottomrule
    \end{tabularx}
\end{frame}

\begin{frame}
    \frametitle{组件 1--3：规则、关系、BDIE}
    \scriptsize
    \begin{block}{1. RuleEngine + 礼单账本（\texttt{rules/rule\_engine.ts}）}
        \textbf{功能}：唯一世界状态写入闸门；所有 cash/ledger/face/trust/prestige 变更经此，写入 append-only 日志（含 before/after/reason）；维护回礼窗口 W4--W10。\\
        \textbf{影响}：保证"内容与决策分离"——LLM 只产出文本与意图，不直接写状态。回礼窗口规则直接支撑 H4（延迟互惠）。
    \end{block}
    \begin{block}{2. 关系图（\texttt{systems/relationship/}）}
        \textbf{功能}：有向边存 trust/intimacy/affection/dislike/reciprocityScore/giftDebt，标记 relationAxis（horizontal/vertical\_up/down）；RelationshipSummarizer 转人话注入 Prompt。\\
        \textbf{影响}：纵向轴是 H3 唯一载体；dislike 是闲话与拒绝目标选择依据；giftDebt 累积是关系深度可计量代理。
    \end{block}
    \begin{block}{3. BDIE（\texttt{cognition/bdi.ts}）}
        \textbf{功能}：BDI + Emotion（mood/stress/energy/arousal）；\texttt{syncRepayIntentionsFromLedger} 把临窗礼债提升为 priority 95--99 的 Intention；\texttt{applyGiftDefaultToBdi} 按 severity 调整情绪与 repair/distance 意图。\\
        \textbf{影响}：使 Agent 决策突破 LLM 默认伦理，做出多样性、务实性决策；揭示情绪—地位—工作—经济双向传播链条。
    \end{block}
\end{frame}

\begin{frame}
    \frametitle{组件 4--6：认知树、PET、LLM 对话}
    \scriptsize
    \begin{block}{4. 认知树（\texttt{cognition/cognitive\_tree.ts}）$\bigstar$}
        \textbf{功能}：四层结构 social $\to$ circle $\to$ relation $\to$ dialogue；对话摘要经 influenceSum 累积，超阈值 15 晋升为 relation 节点，作为"主观关系记忆"注入后续对话 Prompt。\\
        \textbf{影响}：增强 Agent 对公共记忆与个人属性的串联，提供个人认知更新路径。消融 A2 显示：去掉后个人行为变随机，地位差异下礼物单向流动显著增强——证明主观记忆对维持互惠对称性有因果作用。
    \end{block}
    \begin{block}{5. PET 个人事件表（\texttt{systems/event/pet.ts}）}
        \textbf{功能}：每 Agent 维护近期显著事件列表，含 confidence；\texttt{applyPetBeliefs} 同步到 Belief（B $\leftarrow$ PET）。\\
        \textbf{影响}：比全局日志更贴近"个人记得什么"，使话语对齐率 alignmentRate 达到 1.0。
    \end{block}
    \begin{block}{6. LLM 对话系统（\texttt{systems/dialogue/}）}
        \textbf{功能}：多模式（mock/live）对话，受 state\_narrative 约束；LLM 只产文本，不写状态。\\
        \textbf{影响}：话语是过程证据。主跑 344 次对话 0 失败，alignmentRate=1.0，warn=0——证明说与做受同一套约束。
    \end{block}
\end{frame}

\begin{frame}
    \frametitle{组件 7--9：经济、事件、性格}
    \scriptsize
    \begin{block}{7. 经济系统 / 日活力（\texttt{systems/economy/}、\texttt{person/daily\_vitals.ts}）}
        \textbf{功能}：日生活成本 20--90 元、月工资、存款、礼期望膨胀检查；\texttt{dailyLivingCost} 按职业/地位/essentialExpenseRatio 计算。\\
        \textbf{影响}：提供生存约束与压力底噪。消融 A3（无限资金）显示礼期望膨胀、分层扭曲——证明经济约束是分层前提。
    \end{block}
    \begin{block}{8. 事件系统（\texttt{systems/event/}）}
        \textbf{功能}：三类驱动——公共事务（scheduled，如婚礼）、社交驱动（闲话/拜访）、回礼驱动（从礼债生成回礼事件）；\texttt{initiatorProbability}/\texttt{pairScore} 决定"谁和谁互动"。\\
        \textbf{影响}：决定"何时有事"，与 LLM"做什么"解耦——时序真实性的关键，避免同步 slot 退化为 Poisson。
    \end{block}
    \begin{block}{9. 性格 / 人格派生（\texttt{cognition/personality.ts}）}
        \textbf{功能}：从 BigFive 派生 introversion（1--E）、pickiness（0.45(1--A)+0.35N+0.2C）、spitefulness（0.55N+0.45(1--A)）；\texttt{decideRefuse} 用这些特质决定拒绝概率。\\
        \textbf{影响}：把人口统计条件化为异质决策倾向；情绪与个人性格是导致社会分层的重要因素。
    \end{block}
\end{frame}

\section{实验设置}

\begin{frame}
    \frametitle{规模梯度与主跑配置}
    \small
    \begin{columns}
        \begin{column}{0.48\textwidth}
            \textbf{规模梯度}
            \begin{itemize}
                \item 16 人 / 5 天：机制 Demo
                \item 100 人 / 30 天（mock）：对齐短跑，12/12 验收
                \item \textbf{100 人 / 365 天（live）}：主结果
                \item 100 人 / 365 天（live）：seed=42 复跑
            \end{itemize}
        \end{column}
        \begin{column}{0.48\textwidth}
            \textbf{主跑配置}
            \begin{itemize}
                \item Run：\texttt{live100\_365d\_20260725\_1643}
                \item 100 Agents · 365 天 · baseline
                \item DeepSeek · seed=42
                \item 全开：礼单记忆 + 互惠规则 + 场合规范 + BDIE + 认知树 + 经济约束
            \end{itemize}
        \end{column}
    \end{columns}
    \vspace{0.4em}
    \textbf{消融与对照}：A0 全开 / A1 无 BDIE / A2 无认知树 / A3 无限资金；E1 礼单记忆 / E2 互惠规则 / E3 纵向关系 / E4 场合规范。三层指标：宏观 metrics · 微观个体 · 过程（delay/Intention/话语）。
\end{frame}

\section{实验结果}

\begin{frame}
    \frametitle{主跑宏观指标（live100\_365d）}
    \small
    \centering
    \begin{tabularx}{\textwidth}{|l|r|X|c|}
        \toprule
        \textbf{指标} & \textbf{结果} & \textbf{解读} & \textbf{假说} \\
        \midrule
        reciprocity\_rate & 0.90 & 多数还情，互惠中态 & H2 \checkmark \\
        default\_rate & 0.022 & 有违约但不崩溃 & H2 \checkmark \\
        avg\_repay\_delay & $\sim$17 slots & 延迟互惠，非即时会计 & H4 \checkmark \\
        gift\_prestige\_corr & 0.51 & 声望吸引礼物流 & H1 \checkmark \\
        gift\_wealth\_corr & 0.07 & 非纯经济驱动 & H1 \checkmark \\
        Top10 声望礼份额 & 27\% & 中度集中，像分层村庄 & H1 \checkmark \\
        prestige\_gini & 0.16 & 有分层、不极端 & — \\
        dialogues & 344（0 fail） & Live 话语接入过程 & H5 \checkmark \\
        \bottomrule
    \end{tabularx}
\end{frame}

\begin{frame}
    \frametitle{对齐短跑（tune\_align\_v2\_100\_30d）：12/12 通过}
    \scriptsize
    \begin{columns}
        \begin{column}{0.5\textwidth}
            \textbf{情绪与经济带}
            \begin{itemize}
                \item mood $\approx$ 0.55（0.35--0.85）
                \item stress $\approx$ 0.27（0.18--0.72）
                \item energy $\approx$ 0.87（0.45--0.92）
                \item cash 未崩，meanCash=3779
            \end{itemize}
            \textbf{互惠与违约}
            \begin{itemize}
                \item reciprocity $\approx$ 0.86（0.65--1.0）
                \item default $\approx$ 0.025（0.02--0.20）
            \end{itemize}
        \end{column}
        \begin{column}{0.5\textwidth}
            \textbf{过程与话语}
            \begin{itemize}
                \item 拒绝 $\sim$1.3 次/天（有推脱）
                \item gossip=92，dislike$\geq$35 边=7
                \item alignmentRate=1.0，warn=0
                \item 喜事难拒、欠情临窗优先还礼
                \item W4 跨场合还礼日志可查
            \end{itemize}
            \vspace{0.3em}
            \textcolor{blue}{\textbf{意义}}：非全员崩溃、非无压力乌托邦；说与做受同一套约束；人情 $\neq$ 银行即时清算。
        \end{column}
    \end{columns}
\end{frame}

\begin{frame}
    \frametitle{社会分层：职业组对比}
    \tiny
    \centering
    \begin{tabularx}{\textwidth}{|l|c|r|r|r|r|r|r|}
        \toprule
        \textbf{职业组} & \textbf{n} & \textbf{prestige} & \textbf{face} & \textbf{valueIn} & \textbf{valueOut} & \textbf{defaults} & \textbf{degree} \\
        \midrule
        干部/权威 & 17 & \textbf{79.9} & 76.2 & \textbf{1118} & 569 & \textbf{0} & 16.1 \\
        白领/专业 & 25 & 61.4 & 62.2 & 222 & 487 & 0.04 & 15.6 \\
        商业 & 6 & 57.2 & 50.2 & 854 & 332 & 0 & 16.3 \\
        其他 & 8 & 45.7 & 48.2 & 413 & 437 & 0.25 & 17.9 \\
        体力/服务 & 44 & \textbf{42.8} & 42.1 & \textbf{270} & 398 & \textbf{0.14} & 15.2 \\
        \bottomrule
    \end{tabularx}
    \vspace{0.3em}
    \scriptsize
    \textbf{观察}：干部/权威组 prestige=79.9、收礼=1118、违约=0；体力/服务组 prestige=42.8、收礼=270、违约=0.14。\textbf{高地位者收礼多且不违约，低地位者收礼少且违约高}——印证纵向不对称（H3）与声望吸引（H1）。
\end{frame}

\begin{frame}
    \frametitle{社会分层：关键相关与四分位}
    \scriptsize
    \begin{columns}
        \begin{column}{0.5\textwidth}
            \textbf{关键 Pearson 相关}
            \begin{itemize}
                \item face $\leftrightarrow$ reputation = \textbf{0.942}
                \item face $\leftrightarrow$ prestige = 0.913
                \item prestige $\leftrightarrow$ reputation = 0.89
                \item defaults $\leftrightarrow$ giftDebtOut = 0.89
                \item mood $\leftrightarrow$ stress = \textbf{$-$0.817}
                \item income $\leftrightarrow$ face = 0.69
                \item income $\leftrightarrow$ prestige = 0.688
                \item \textbf{prestige $\leftrightarrow$ valueIn = 0.575}
                \item wealth $\leftrightarrow$ valueIn = 0.109
            \end{itemize}
        \end{column}
        \begin{column}{0.5\textwidth}
            \textbf{四分位对比}
            \begin{itemize}
                \item 高 prestige Q4 vs Q1：收礼差 \textbf{644}，face 差 38.3，degree 差 0.7
                \item 高 wealth Q4 vs Q1：prestige 差仅 0.8，收礼差 200
            \end{itemize}
            \vspace{0.3em}
            \textcolor{red}{\textbf{结论}}：\textbf{"位 $\to$ 礼"而非"钱 $\to$ 礼"}；情绪与地位双向耦合；脸面-声望共变而非脱钩。
        \end{column}
    \end{columns}
\end{frame}

\begin{frame}
    \frametitle{消融：证明"组件独特作用"}
    \small
    \begin{tabularx}{\textwidth}{|l|X|}
        \toprule
        \textbf{消融} & \textbf{预期（若组件真有用）} \\
        \midrule
        \textbf{A2 无认知树} & 对话语境与主观关系史变弱；长期互动质量下降；地位差异下礼物单向流动显著增强 \\
        \midrule
        A1 无 BDIE & 社交发起/拒绝/回礼紧迫度"变平"；决策同质化 \\
        \midrule
        A3 无限资金 & 礼期望膨胀，经济约束假说被破坏，分层扭曲 \\
        \midrule
        E2 互惠规则关 & 违约 $\uparrow$、互惠 $\downarrow$ \\
        \midrule
        E3 仅横向 & 纵向不对称失效，进贡不再单向 \\
        \midrule
        E4 无场合规范 & "喜事难拒"消失，拒绝率异常上升 \\
        \bottomrule
    \end{tabularx}
    \vspace{0.3em}
    \scriptsize
    \textbf{话术}：消融不是凑材料，而是回答——"多智能体多出来的那些模块，是不是真的在干活？"
\end{frame}

\section{结论向实验结果的解释}

\begin{frame}
    \frametitle{结论 1：Multiagent 拟真较好，决策能力相似}
    \small
    \begin{itemize}
        \item \textbf{证据}：reciprocity=0.90、default=0.022、delay$\approx$17 slots 落在民族志"互惠中态"区间；话语 alignmentRate=1.0 说明决策与话语受同一套约束。
        \item \textbf{机制}：BDIE 把"手头紧/高压/欠情临窗"变成可计算驱动——\texttt{syncRepayIntentionsFromLedger} 把临窗礼债提升为 priority 95--99 的 Intention，使 Agent 决策突破 LLM 默认伦理，做出多样性、务实性选择。
        \item \textbf{多层次有效}：个人层（情绪带/性格）、集体层（7272 events/92 gossip）、社会层（职业分层 prestige 79.9 vs 42.8）。
        \item \textbf{可解释}：所有状态变更经 RuleEngine 写入 append-only 日志（含 before/after/reason）；BDIE 轨迹、认知树晋升、话语对齐率均可回放。
        \item \textbf{可复现}：seed=42 复跑所有指标方向一致、量级同阶；mock 对齐短跑 12/12 通过。
    \end{itemize}
\end{frame}

\begin{frame}
    \frametitle{结论 2：认知树与 BDIE 的独特作用}
    \small
    \begin{block}{认知树：主观关系记忆的因果作用}
        \textbf{功能证据}：对话摘要经 influenceSum 累积晋升为 relation 节点，注入后续对话 Prompt，使 alignmentRate=1.0。\\
        \textbf{消融 A2 证据}：去掉认知树后，个人行为变随机，且在地位差异情况下\textbf{礼物单向流动的特征显著增强}——证明主观记忆对维持互惠对称性有因果作用。
    \end{block}
    \begin{block}{BDIE：突破 LLM 伦理，揭示双向传播}
        \textbf{证据}：mood $\leftrightarrow$ stress = $-$0.817 揭示情绪双向传播；income $\leftrightarrow$ face = 0.69、income $\leftrightarrow$ prestige = 0.688 揭示"情绪—社会地位—工作内容—经济状态"\textbf{双向链条}，而非单向。\\
        \textbf{结论}：情绪和个人性格是导致社会分层的重要因素，进而导致社会分层。
    \end{block}
\end{frame}

\begin{frame}
    \frametitle{结论 3：关系网络与集体事务效率}
    \small
    \begin{block}{关系网络：频率主、礼债时长次}
        \textbf{证据}：prestige $\leftrightarrow$ valueIn = 0.575（频率/价值驱动），wealth $\leftrightarrow$ valueIn = 0.109（财富不驱动）；defaults $\leftrightarrow$ giftDebtOut = 0.89（礼债累积与违约同向）。\\
        \textbf{机制}：giftDebt 持续时间越长，trust/intimacy/reciprocityScore 越高（关系越深），但违约风险也越高。
    \end{block}
    \begin{block}{集体事务效率：影响因素与权重}
        \textbf{证据}：7272 events 完成，干部组 defaults=0、degree=16.1（高网络度数 + 高地位 $\to$ 集体事务解决效率高）；体力组 defaults=0.14、degree=15.2（低地位 $\to$ 违约率高）。\\
        \textbf{影响程度排序}：\textbf{社会地位 $>$ 关系网络 $>$ 经济状况 $>$ 性格}。
    \end{block}
\end{frame}

\section{可说明结论的更多实验变量}

\begin{frame}
    \frametitle{已具备的实验变量}
    \tiny
    \begin{tabularx}{\textwidth}{|l|l|X|}
        \toprule
        \textbf{变量} & \textbf{脚本} & \textbf{说明的结论} \\
        \midrule
        A2 无认知树 & run\_ablation.ts & 认知树对主观关系记忆的因果作用；去掉后单向流动增强 \\
        \midrule
        A1 无 BDIE & run\_ablation.ts & BDIE 对决策多样性的贡献；去掉后社交/回礼变"平均人" \\
        \midrule
        A3 无限资金 & run\_ablation.ts & 经济约束对分层的前提作用；去掉后礼期望膨胀 \\
        \midrule
        E2 互惠规则关 & run\_contrast.ts & 互惠规则对违约/互惠率的因果作用 \\
        \midrule
        E3 仅横向 & run\_contrast.ts & 纵向轴对单向流动的因果作用 \\
        \midrule
        E4 无场合规范 & run\_contrast.ts & 场合规范对"喜事难拒"的作用 \\
        \midrule
        seed=42 复跑 & run\_experiment.ts & 可复现性 \\
        \midrule
        mock vs live & experiment\_presets.ts & LLM 话语对过程的贡献 \\
        \bottomrule
    \end{tabularx}
\end{frame}

\begin{frame}
    \frametitle{建议补充的实验变量}
    \tiny
    \begin{tabularx}{\textwidth}{|l|X|X|}
        \toprule
        \textbf{变量} & \textbf{设计} & \textbf{说明的结论} \\
        \midrule
        V1 关系网络密度梯度 & 初始 edge 密度 0.1 / 0.25 / 0.5 & 关系网络对集体事务效率的影响程度 \\
        \midrule
        V2 性格异质性梯度 & 全员中性 vs BigFive 异质 vs 极端异质 & 性格对决策多样性与分层的作用 \\
        \midrule
        V3 社会地位初始分布 & 均匀 vs 偏态（少数高地位） & 地位差异对礼物单向流动的放大效应 \\
        \midrule
        V4 经济状况梯度 & 贫困 / 中等 / 富裕村庄 & 经济状况对礼期望与分层的影响程度 \\
        \midrule
        V5 回礼窗口长度 & W2 / W4 / W10 / W20 & 时间性对延迟互惠与违约率的作用 \\
        \midrule
        V6 礼债时长 × 关系深度 & 跟踪 giftDebt 时长与 trust/intimacy 时序相关 & 礼债持续时间作为关系深度次要因素的量化 \\
        \midrule
        V7 LLM 模型替换 & DeepSeek / Qwen / GPT & 模型敏感性与结论稳健性 \\
        \midrule
        V8 温度敏感性 & T=0.2 / 0.7 / 1.0 & 决策确定性与多样性的权衡 \\
        \midrule
        V9 集体事务难度梯度 & 简单（婚礼）/ 中等（公共工程）/ 困难（纠纷调解） & 关系/性格/地位/经济对解决效率的相对权重 \\
        \midrule
        V10 反事实：切断纵向轴 & E3 强化版，跟踪 face/trust 轨迹 & 纵向不对称对分层的因果作用 \\
        \bottomrule
    \end{tabularx}
\end{frame}

\begin{frame}
    \frametitle{实验变量与结论的映射矩阵}
    \scriptsize
    \begin{itemize}
        \item \textbf{Multiagent 拟真较好}：A1, mock vs live, V7, V8
        \item \textbf{多层次属性有效模拟}：V1, V2, V3, V4
        \item \textbf{可解释性}：所有 A/E（日志可回放）
        \item \textbf{可复现性}：seed=42 复跑, V7, V8
        \item \textbf{认知树作用}：A2, V6
        \item \textbf{BDIE 突破 LLM 伦理}：A1, V2, V8
        \item \textbf{情绪-地位-工作-经济双向链条}：A1, V4, 分层相关分析
        \item \textbf{关系网络：频率主、礼债时长次}：V1, V6, E3
        \item \textbf{集体事务效率影响因素}：V1, V3, V4, V9
    \end{itemize}
\end{frame}

\section{局限与结论}

\begin{frame}
    \frametitle{局限（坦诚反而加分）}
    \small
    \begin{enumerate}
        \item \textbf{宏观-微观解离风险}：群体指标像民族志，不证明每个 Agent 的神经/心理过程像人。
        \item \textbf{人类数据对齐仍浅}：尚未做问卷边际 / 分布级 KS 对照。
        \item \textbf{时序简化}：同步 slot $\neq$ 真实 burstiness。
        \item \textbf{模型敏感性}：主结果单一 live 模型（DeepSeek）。
        \item \textbf{伦理边界}：机制探索与教学，不替代田野因果识别。
    \end{enumerate}
\end{frame}

\begin{frame}
    \frametitle{结论：五个"可取之处"}
    \small
    \begin{enumerate}
        \item \textbf{拟真可检验}：100 人村庄复现声望吸引的礼物集中（$r=0.51$）与延迟互惠中态（reciprocity=0.90, delay$\approx$17 slots），短跑对齐带 12/12 通过。
        \item \textbf{组件有独特作用}：规则负责规范承诺，BDIE 负责状态驱动，\textbf{认知树负责主观关系记忆}，LLM 负责可回放话语——彼此不可替代，可用消融追问。
        \item \textbf{方法可取}：多智能体社会拟真把人情社会变成"可开关机制 + 三层指标 + 可回放过程"的实验装置，服务于理论筛选与教学演示。
        \item \textbf{情绪与性格导致分层}：情绪—社会地位—工作内容—经济状态的双向传播链条是分层的重要机制，而非单向链条。
        \item \textbf{关系网络量化}：礼物流动频率是关系深度的主要因素，礼债持续时间是次要因素。
    \end{enumerate}
\end{frame}

\begin{frame}
    \frametitle{下一步}
    \small
    \begin{itemize}
        \item 消融对比幻灯（尤其 A2 认知树定量对比）
        \item 与民族志/问卷指标做分布级对照（KS 检验）
        \item 模型敏感性（另一 LLM / 温度）
        \item 反事实：切断纵向轴或关闭场合规范
        \item 过程指标仪表盘：BDI 轨迹、认知树晋升、话语对齐率纳入例行报告
        \item 补充建议变量（V1--V10），量化各因素对集体事务效率的相对权重
    \end{itemize}
\end{frame}

\begin{frame}
    \begin{center}
        {\Huge 感谢聆听！}
    \end{center}
\end{frame}

\end{document}

